% TOP_PROBLEM: {"id":30007610,"problem_number":"LOCAL-30007610","title":"Platform credit and inclusion","category_id":21,"category":{"id":21,"name":"economics","display_name":"Economics","description":"Open economic theory statements, published research agendas, and proposed scoped specifications; question provenance and historical priority are qualified per record.","slug":"economics","order_index":21,"created_at":"2026-10-08T00:00:00Z"},"status":"research_question","external_url":null,"legacy_ids":[],"published":false,"statement":"In the explicitly specified setting: When does platform-based lending expand useful credit access, and when does data concentration create market power or excessive risk? The mathematical statement above fixes the target, assumptions and scope of this catalogue formulation.","economics":{"original_id":99,"field":"Banking, credit and financial stability","kind":"design","formalization_basis":"adapted_research_question","source_status":"explicit_open","priority_status":"earliest_found","priority_note":"This is the earliest explicit statement verified in this audit; first priority for the full catalogue formulation is not established.","earliest_verified_reference":{"authors":["Karen Croxson","Jon Frost","Leonardo Gambacorta","Tommaso Valletti"],"title":"Platform-based business models and financial inclusion","year":2022,"url":"https://www.bis.org/publications/working-paper-986-platform-based-business-models-and-financial-inclusion.pdf","doi":null,"locator":"Section 1, Introduction, printed p.2 (PDF p.4), second paragraph","evidence_summary":"Explicitly identifies inclusion, competition, stability and privacy tradeoffs."},"current_reference":{"authors":["Karen Croxson","Jon Frost","Leonardo Gambacorta","Tommaso Valletti"],"title":"Platform-based business models and financial inclusion","year":2022,"url":"https://www.bis.org/publications/working-paper-986-platform-based-business-models-and-financial-inclusion.pdf","doi":null,"locator":"Section 1, Introduction, printed p.2 (PDF p.4), second paragraph","evidence_summary":"Explicitly identifies inclusion, competition, stability and privacy tradeoffs."},"formalization_note":"Adapts an explicitly verified research-agenda passage. Mathematical objects, interventions, objectives and scope are proposed here; existing model-specific solutions are not relabeled unsolved."},"statusQualification":"Published question or research agenda verified at the cited locator. The mathematical specification is adapted for this catalogue and includes additional assumptions; source publication does not establish first-ever priority or certify this exact model remains unsolved. Adapts an explicitly verified research-agenda passage. Mathematical objects, interventions, objectives and scope are proposed here; existing model-specific solutions are not relabeled unsolved.","statusReviewedAt":"2026-10-08"}
% FIRST_POSTED: 2026-10-08
% ENTRY_KIND: statement_only
% !TeX program = lualatex
\documentclass[11pt]{article}
\usepackage{fontspec}
\usepackage[a4paper,margin=25mm]{geometry}
\usepackage{amsmath,amssymb,mathtools}
\usepackage{xurl}
\usepackage[hidelinks]{hyperref}
\newcommand{\sref}[2]{\hyperlink{source:#1:#2}{[S#2]}}
\setlength{\emergencystretch}{3em}
\begin{document}
% problemId:30007610
\section*{Platform credit and inclusion}\label{30007610}
\subsection{Definitions and mathematical statement}
Consider a two-sided lending platform linking borrowers \(i\) to lenders, using borrower data to screen risk. Let \(n\) measure the platform's network size, \(d\) its data access, and \(a\) a policy specifying interoperability, data portability, privacy restrictions and entry conditions. Equilibrium loan prices, approval probabilities and default losses depend on \((n,d,a)\), with platform entry and borrower participation endogenous. Define inclusion as access to loans that increase borrowers' expected utility net of repayment, not merely observed approval. Let \(W(a)\) combine borrower and lender surplus, real screening costs, privacy losses and systemic externalities. The task is to characterize the welfare frontier
\[\{(\text{useful credit access},\text{market power},\text{privacy loss},\text{default risk})(a):a\in A\}.\]
Here \(A\) is the stated feasible policy set and each component has a declared population and measurement rule. Identify when additional data improves screening and when it mainly strengthens incumbent advantage. Compare policies holding primitives fixed while allowing equilibrium adaptation. State the exclusions needed to identify data and network effects separately; selection into platforms can otherwise mimic improved lending. Determine the conditions under which portability raises inclusion without worsening risk or privacy outcomes.

\par\textbf{Formulation and scope.} Adapts an explicitly verified research-agenda passage. Mathematical objects, interventions, objectives and scope are proposed here; existing model-specific solutions are not relabeled unsolved.

\par\textbf{Open-question provenance.} An explicit question or research agenda was verified at the source locator. This does not by itself certify that every later version remains unresolved \sref{30007610}{1}.

\par\textbf{Historical priority.} This is the earliest explicit statement verified in this audit; first priority for the full catalogue formulation is not established.

\subsection{Short English statement}
In the explicitly specified setting: When does platform-based lending expand useful credit access, and when does data concentration create market power or excessive risk? The mathematical statement above fixes the target, assumptions and scope of this catalogue formulation.

\subsection{Sources}
\begin{itemize}\raggedright
\item[\textnormal{[S1]}]\hypertarget{source:30007610:1}{} Karen Croxson, Jon Frost, Leonardo Gambacorta, Tommaso Valletti. 2022. Platform-based business models and financial inclusion. Section 1, Introduction, printed p.2 (PDF p.4), second paragraph.
\url{https://www.bis.org/publications/working-paper-986-platform-based-business-models-and-financial-inclusion.pdf}.
{\small\textit{Earliest verified question/agenda reference; historical priority qualified below. Explicitly identifies inclusion, competition, stability and privacy tradeoffs.}}
\end{itemize}
\end{document}
